\documentclass[11pt,a4paper]{article}
\usepackage[T1]{fontenc}
\usepackage{lmodern,amsmath,amssymb,amsthm,booktabs,microtype,hyperref,array}
\pdfmapfile{+lm.map}
\usepackage[margin=26mm]{geometry}
\setlength{\emergencystretch}{2em}
\newcommand{\Q}{\mathbb Q}\newcommand{\Z}{\mathbb Z}
\newcommand{\Tr}{\operatorname{Tr}}\newcommand{\N}{\operatorname{N}}
\newtheorem{proposition}{Proposition}\newtheorem{lemma}[proposition]{Lemma}
\title{Ancillary audit: intrinsic slices, selector density, and local normalization}
\author{Z diophantine in $\Q$ project}
\date{19 November 2026}
\begin{document}\maketitle
\begin{abstract}
We audit the attached realization, selector, and point-restriction results.
The two explicit universal-specialization counterexamples are valid, and
an independent exact derivative certificate verifies the rank-one theorem
used by the zero-ordinate obstruction. Their scopes must be preserved:
neither counterexample is known to be an actual detector input that must
occur for an integer. We also prove a local trace lemma and verify a
birational intrinsic-pole normalization used in the improved ordinary
eleven-variable source.
\end{abstract}

\paragraph{Source credit.}
The supplied constructions and notes audited here are by Gabe Conant with
Astra~6.2, in the research bundles
\emph{H10-quartic-constructions} and
\emph{H10-eleven-variable-constructions}.
The audit, corrections, and further improvements stated in this paper are
subsequent project work. The original source files and their provenance
remain separately preserved.

\section{Audit verdict and scope}
The following attached statements withstand the audit. ``Prepared'' means
$0<d<1$, a square reduced denominator, $d\in7\Q_2^{\times2}$, and
$v_2(d)\ge4$. The finite checks supplement the indicated proofs; they do
not replace their quantifiers.

\begin{center}\small
\begin{tabular}{p{.28\linewidth}p{.65\linewidth}}
\toprule
Attached result & Audited scope\\\midrule
Every diagonal prepared fiber is everywhere locally soluble &
Valid, with nonzero affine ordinate at every completion. Rational weak
approximation of the abscissa supplies local squares at a prescribed finite
set, not a global square.\\[3pt]
$d=112/289$ diagonal failure &
Valid global failure with all completions soluble. It disproves uniform
diagonal realization on all prepared inputs. It is not identified as a
forced actual detector value.\\[3pt]
Infinite soluble diagonal families &
Valid. One has fixed global class $7$ and nonzero ordinate; the other has
zero ordinate and escapes every finite set of classes for $d$ and $1-d$.
Neither proves completeness of a diagonal defining source.\\[3pt]
Selector-density obstruction &
Valid under the inherited supplier and density premises for the complete
set of all old prepared tuples. It is not a general lower bound of eleven
or thirteen variables.\\[3pt]
Fixed-twist slice finiteness &
Valid for fixed twist and a plane constraint without horizontal or vertical
geometric components. The twist may vary in a proposed new construction.\\[3pt]
$d=112/225$ zero-ordinate obstruction &
Valid. The required elliptic rank-one input is unconditionally certified
using the standard analytic formula and Gross--Zagier--Kolyvagin. Again this
is a universal prepared-input obstruction, not actual detector alignment.\\[3pt]
Fixed proper point restrictions &
Valid as an obstruction to arbitrarily fresh prime selection on a fixed
proper algebraic locus. It does not exclude using a finite set of primes
or input-dependent loci.\\\bottomrule
\end{tabular}
\end{center}

\section{The diagonal Hasse-principle failure}
The diagonal curve is
\[
 C_d:\quad Y^2=d\big[(d+v^2)^2v^4-(d-v^2)^2(1-d)^2\big].
\]
Writing $r=v^2,h=1-d$, its bracket factors into
\[
 (r^2+(d-h)r+dh)(r^2+(d+h)r-dh).
\]
The two factors are coprime; their discriminants are
$d^2-6dh+h^2$ and $d^2+6dh+h^2$. Neither can vanish for rational
$d/h$, and their constant terms are nonzero. This proves that the eight
roots in $v$ are distinct and $C_d$ has genus three.

The local proof in the attachment is sound, including its exceptional
cases. Where $d$ is square, sufficiently large $v$ works; where $-d$ is
square, $v^2=-d$ works. The latter includes $2$. At odd primes with
$v_p(d)>0$, the bracket has a simple root at $v=1$ modulo $p$; lift it
and perturb the smooth zero-ordinate point. In the remaining nonsquare
cases $p\equiv1\pmod4$. After normalization the equation is
\[
 y^2=D\big[(D+x^2)^2x^4-(D-x^2)^2h^2\big]
\]
with $D$ a nonsquare unit and $h\equiv1-D$ or $-D$. If the two quadratics
are separable, the genus-three bound gives a point for $p\ge37$, and the
two points at infinity are not rational. If one quadratic is a square,
the remaining genus-one curve has enough points for $p\ge13$ to avoid
the at most four excluded points. The discriminants differ by $12Dh$,
so they cannot both vanish. An independent symbolic derivative and
finite-field enumeration checks all $60$ remaining cases for
$p=5,13,17,29$.

For $d=112/289$ and primitive $(a,b)$, let $r=a^2,s=b^2$ and
\[
 U_+=83521r^2-18785rs+19824s^2,\qquad
 U_-=83521r^2+83521rs-19824s^2.
\]
Their product must have square class $7$. The elementary half-sum and
half-difference argument restricts their common prime support to
$S=\{2,3,7,17,59\}$. Therefore, including the cases where one factor
vanishes, some signed squarefree $S$-supported class $k$ makes
$kU_+$ and $jU_-$ squares, where $j$ represents $7k$.

Our independent checker enumerates \emph{all} primitive residue pairs,
instead of importing the attachment's projective-chart enumeration:
\[
\begin{array}{c|r|r}
\text{modulus}&\text{primitive pairs checked}&\text{classes surviving}\\\hline
256&49152&32\\27&648&16\\49&2352&4\\41&1680&0
\end{array}
\]
The four classes before the last step are $(7,1),(14,2),(-413,-59)$,
and $(-826,-118)$. The points at infinity are nonrational because $d$
is nonsquare. Thus the smooth projective curve has no rational point.
The attachment's generalization to fixed slopes whose prime factors
are $1$ modulo $256\cdot27\cdot49\cdot41=13886208$ follows by the
same finite residue patterns and is valid.

The distinct elementary valuation obstruction to $s=cv^m$, $m\ge2$,
is also valid. At a sufficiently large prime $p\equiv13\pmod{28}$,
take $d=112/p^2$. The two terms force $v_p(s)=-1$, impossible when
$v_p(c)=0$ and $m\mid v_p(s)$. For a rational function $c(d)$ this
gives the necessary condition $m\mid 2\deg_\infty(c)-1$. It remains a
condition on this class of specializations, not all possible couplings.

The two infinite globally soluble families in the attachment use
nontorsion seeds on elliptic quartics. Their displayed maps to integral
Weierstrass models preserve a denominator at $67$ in the first case;
in the second $x(2(8,24))=49/36$. Nagell--Lutz supplies nontorsion.
Local logarithms, the Chinese remainder theorem, and real circle density
justify the simultaneous open conditions. In the second family, the two
double covers adjoining a prescribed class of $d$ or $1-d$ ramify at
simple zeros, have genus at least two, and Faltings gives the asserted
escape from any finite set of square classes. These arguments do not
restore the lost two-dimensional pole-parameter density on a fixed
genus-three diagonal fiber.

\section{The zero-ordinate obstruction and its rank certificate}
Take $d=112/225=7(4/15)^2$ and $1-d=113/225$. For a putative zero
ordinate in the intrinsic pole, set
\[
 e=\frac{d-v^2}{d+v^2},\qquad p=\frac{2dv}{d+v^2}.
\]
Then $p^2=d(1-e^2)$ and $s^2=(1-d)|e|$. The exceptional $e=0$ is
impossible because $d$ is nonsquare; $s=0$ would force the same
exception. Put $N=791=7\cdot113$ and
\[
 a=\frac{113}{4s},\qquad
 x=\frac{N(1-d)}{s^2}=\frac{N}{|e|}=da^2,\qquad y=axp.
\]
Then $y^2=x^3(1-e^2)=x^3-N^2x$, giving a rational
point on $E_N$ with square class $x=7\Q^{\times2}$.

Here is a complete certificate route for ruling it out. The conductor is
$32N^2=20021792$ and the root number is $-1$, since $N\equiv7\pmod8$.
These standard formulas for congruent-number curves are explicitly
recorded in the primary sources \cite{Stein,Elkies}. The analytic formula
implemented in Sage's documented \texttt{deriv\_at1} routine \cite{Sage} is
\[
 L'(E_N,1)=2\sum_{n\ge1}\frac{a_n}{n}E_1(cn),\qquad
 c=\frac{2\pi}{\sqrt{32N^2}},\quad
 E_1(cn)=\int_n^\infty\frac{e^{-ct}}{t}\,dt.
\]
The supplied attachment uses a mesh of $1/16$ and cutoff $12000$.
We independently use a unit mesh and cutoff $7000$, direct
Legendre-symbol counts for every good prime, and prime-power recurrences
followed by multiplicativity. There is no call to a numerical rank routine.

For reproducibility, Machin's identity
$\pi=16\arctan(1/5)-4\arctan(1/239)$ and alternating-series remainder
bounds give exact rational bounds for $c$. Degree-seven and degree-six
alternating Taylor bounds enclose $e^{-c}$. All powers and sums are
rounded outward with integer scale $10^{30}$. The decreasing integrand
is enclosed by its right and left unit-mesh sums. The omitted mesh is
bounded geometrically. Finally $|a_n|\le\tau(n)\sqrt n\le2n$ gives
\[
 2\sum_{n>K}\frac{|a_n|}{n}E_1(cn)
 \le\frac{4e^{-c(K+1)}}{c(K+1)(1-e^{-c})}.
\]
The exact resulting interval is
\[
 \frac{6489867793040175244734154791093}{10^{30}}
 \le L'(E_{791},1)\le
 \frac{8941913002135842990365653860323}{10^{30}}.
\]
It is strictly positive. Functional-equation parity therefore gives
analytic rank exactly one. Modularity and Gross--Zagier--Kolyvagin give
algebraic rank exactly one; this implication is unconditional and is
stated explicitly in \cite{Venerucci}. No Birch--Swinnerton-Dyer conjecture
is being assumed.

Reduction at $3$ and $19$, with group orders $4$ and $20$, shows that
the torsion is the full rational two-torsion. The point
$G=(-784,-2940)$ lies on $E_{791}$. For the Kummer map $x+791$, the
torsion image is $\{1,791,1582,2\}$ but $x(G)+791=7$ is outside it.
Thus $G$ is an odd multiple of a free generator modulo torsion. Under
the Kummer map $x$, the torsion image is
$\{1,-1,791,-791\}$ and $x(G)$ has class $-1$. Rank one then forces
the entire image of $x$ to lie in that four-element group. The class $7$
is absent, as required.

The related quartic $y^2=113(112^2-z^4)$ is everywhere locally soluble:
at good primes use a smooth genus-one point and local perturbation;
at $2,7$ use the local squareness of $113$; at $113$ use
$z^2=112$ and perturb the rational zero ordinate. This does not conflict
with the preceding global obstruction.

\section{Selector and point restrictions}
The density proof correctly starts from the upstream CM-independence
lemma for the two conjugate generators. It uses the identity component
of their compact real orbit closure, whose nonempty real open subsets
are Zariski dense. The multiplication graphs
$(n,[n]P,H,[n]H)$ are dense as $n$ varies: divide a putative polynomial
relation by its largest power of $n$, approximate arbitrary triples in
the real torus using the $1/n$ preimage mesh, and descend on that power.
The supplier's reference progressions and unbounded slopes, followed by
pole-parameter density, give density of \emph{all} old prepared base
tuples. Taking products under the sign changes handles square-root
abscissa charts.

Consequently any fixed rational substitution for the selector, or finite
collection of polynomial conditions, which vanishes on every old base
tuple must vanish identically after denominator clearing. This excludes
retaining every old tuple with the same free-coordinate interface. It
does not exclude a new proper slice with at least one suitable point
over each integer.

For a fixed intrinsic twist and nonsquare $1-d>0$, the two elliptic
factors $E_n$ and $E_{n(1-d)}$ are not $\Q$-isogenous. Indeed their
geometric homomorphisms are defined over
$\Q(i,\sqrt{1-d})$, and the automorphism fixing $i$ but negating
$\sqrt{1-d}$ negates every such homomorphism. A plane constraint with
no horizontal or vertical component pulls back to proper curves.
Mordell--Lang would make an infinite component a translate of an
elliptic subgroup; a subgroup projecting nontrivially to both factors
would be an isogeny, and factor subgroups are excluded by the no-line
hypothesis. Thus the fixed-twist finiteness assertion is valid.

The affine index-restriction proof also correctly allows half-integral
auxiliary indices. Doubling kills two-torsion; successful local tests
then force one of two explicitly fixed nonidentity multiples of $P$ to
reduce to the origin. Such a point has only finitely many reduction
primes. Mordell--Lang extends this to a fixed finite menu of proper
algebraic point loci. Its conclusion is about fresh-prime supply, and
should not be promoted to a general variable-count obstruction.

Two further restricted-family statements also have the claimed scopes.
If freezing all other data of a genuine witness leaves pole inputs
$b_i(X)=d_i r_i(X)^2$ with fixed classes and leaves only open residual
conditions on $X$, weak approximation preserves those conditions at the
finite exceptional places and the real place while giving $X$ a
denominator at a new prime. Outside the exceptional set the pole
valuations remain even. The prepared converse and finite-character
avoidance then produce a noninteger input. This excludes precisely such
fixed-class perturbation schemes; additional genuine equations in $X$
or varying classes are outside its hypotheses.

The vertical/horizontal three-fiber lemma is likewise valid: if the
divisors of $f$ and $1-f$ on a product of smooth projective curves have
only vertical and horizontal components, write
$f=u(x)v(y)$ and $1-f=r(x)s(y)$. The tensor identity
$1=u\otimes v+r\otimes s$ has rank one. One pair of factors must be
linearly dependent; the identity then forces $f$ to depend on only one
coordinate. This statement constrains that divisor-support pattern.

Finally the pair $d_0=112/13^2,d_1=112/13^4$ disproves a shared diagonal
parameter theorem for arbitrary prepared pairs: at $13$, both $d_i$
and $-d_i$ are nonsquares and the two radicand terms can avoid
nonsquare dominance only when $v_{13}(v)=-1$ and $-2$, respectively.
No single $v$ meets both conditions. Actual alignment of detector pairs
has further structure and is not excluded by this example.

\section{A rigorous trace lemma for the new intrinsic normalization}
Let $K=\Q(\sqrt2)$, $D=11-6\sqrt2$, and
\[
 P=\big((57-28\sqrt2)/8,(-660+417\sqrt2)/32\big),\qquad
 T(n)=\Tr x([n]P).
\]
\begin{lemma}
For every nonzero integer $n$,
\[
 n^2T(n)-57/4\in16\Z_2,\qquad
 v_2(T(n))=-2-2v_2(n),\qquad T(n)\in\Q_2^{\times2}.
\]
Consequently, for integral $\beta$ and any integer $e\ge2$,
$\beta^e+T(n)^e$ is a nonzero square in $\Q_2$.
\end{lemma}
\begin{proof}
Put $\pi=\sqrt2$ and $\mathcal O=\Z_2[\sqrt2]$, with $v_\pi(\pi)=1$.
The formal parameter $t=-x/y$ at $P$ has valuation three, while
$v_\pi(x(P))=-6$. The integral formal group commutes with the CM
automorphism $t\mapsto it$. Thus for odd $m$ its multiplication series
has only powers $1$ modulo four:
\[
 [m](t)=mt+O(t^5),\qquad x(t)=t^{-2}+O(t^2).
\]
All coefficients here are integral in the relevant valuation ring.
Since $m$ is a unit, this proves
$m^2x([m]P)-x(P)\in\pi^6\mathcal O$.

For any formal point $Q$, the exact duplication identity is
\[
 4x([2]Q)-x(Q)
  =\frac{D(3x(Q)^2+D)}{x(Q)(x(Q)^2-D)}.
\]
If $v_\pi(x(Q))<0$, its right side has valuation $-v_\pi(x(Q))$;
also $v_\pi(x([2]Q))=v_\pi(x(Q))-4$.
For $n=2^rm$ with $m$ odd, telescope the duplication identities after
multiplication by $4^jm^2$. Every added term lies in
$\pi^6\mathcal O$. Hence $n^2x([n]P)-x(P)\in\pi^6\mathcal O$.
The trace of $\pi^6\mathcal O=8\mathcal O$ is contained in $16\Z_2$,
and $\Tr x(P)=57/4$. This proves the first assertion. Multiplying by
four gives $4n^2T(n)\equiv57\equiv1\pmod8$, proving both the valuation
and square-class assertions. Finally $v_2(\beta/T(n))\ge2$, so
\[
 \frac{\beta^e+T(n)^e}{T(n)^e}
      =1+(\beta/T(n))^e\in1+16\Z_2\subset\Q_2^{\times2}.
\]
Since $T(n)$ itself is square, this works for odd as well as even $e$.
\end{proof}

This final assertion concerns the local square class at $2$. A global
square reduced denominator is automatic for even $e$, since at an odd
denominator prime of $T$ its valuation in $\beta^e+T^e$ is $e\,v_p(T)$.
For odd $e$ that parity statement is not asserted. In particular
$M=\beta^4+T^4$ and every positive integral power of $M$ have the needed
square denominator, and $M$ itself has the required local square class.
The separate odd-$e$ weak-$abc$ generalization absorbs the fixed scale
denominators in its finite preparation and does not need the stronger
global claim.

Finally suppose the intrinsic input has the exact factorization
$\eta=dt^2$, with $t\ne0$. In the original radicand substitute
$v_{\rm old}=tv$ and $Y_{\rm old}=t^3Y$. Cancelling $t^6$ gives
\[
 Y^2=d(d+v^2)^2s^4-d(d-v^2)^2(1-\eta)^2.
\]
This is a birational change of the two pole parameters and ordinate, with
an explicit inverse. It preserves rational existence, zero-ordinate
cases, and density for avoiding finitely many proper conditions. The
independent checker verifies both this polynomial identity and the exact
duplication identity above.

\begin{thebibliography}{9}
\bibitem{Attachment} Gabe Conant with Astra~6.2, supplied research package,
\emph{H10-eleven-variable-constructions},
\texttt{counts/realization/}, \texttt{counts/selector/}, and
\texttt{counts/points/point-restrictions.md}, 18 November 2026.
\bibitem{Project} Project source at
\texttt{d3a8c223716a0fc0dbf5af4eb7998d640c21048d}:
weighted trace-height v27 (CM independence), pole density v14,
intrinsic pole v56/v57, derivative detector v63, full auxiliary v71,
and ordinary scale v75. These are inherited arithmetic inputs where
explicitly indicated, not claims proved by the present finite checkers.
\bibitem{Stein} W. Stein, \emph{The Congruent Number Problem}, SIMUW notes.
\url{https://www.williamstein.org/simuw06/notes/notes.pdf}.
Conductor and Euler-coefficient formulas for the congruent-number family.
\bibitem{Elkies} N. Elkies, \emph{Elliptic curves $ny^2=x^3-x$ of rank 3,
$|n|\equiv7\pmod8$, $|n|<10^6$}.
\url{https://people.math.harvard.edu/~elkies/cong_r3_7a.html}.
The functional-equation sign and its residue-class formula.
\bibitem{Sage} SageMath reference manual, \emph{L-series of elliptic curves},
\texttt{deriv\_at1}.
\url{https://doc.sagemath.org/html/en/reference/arithmetic_curves/sage/schemes/elliptic_curves/lseries_ell.html}.
\bibitem{Venerucci} R. Venerucci, \emph{On the p-converse of the
Kolyvagin--Gross--Zagier theorem}, introduction, the forward
Gross--Zagier--Kolyvagin implication for analytic rank at most one.
\url{https://arxiv.org/abs/1407.1816}.
The additional hypotheses for converse theorems in that article are not
being used here.
\end{thebibliography}
\end{document}
